% Chapter 6 — Conclusions and Future Work
% Plan mapping: PLANO §6.1 (row 6): demonstrated contributions, limitations,
% future work. Indicative target: 1,500–2,000 words. Contribution and RQ
% statements are finalised only after Chapter 5 is populated from data-v1.

\chapter{Conclusions and Future Work}
\label{ch:conclusions}

\section{Summary}
\label{sec:concl-summary}

\noindent This dissertation designed, implemented and evaluated a reproducible
ARM64 edge gateway for digital twins of wearable devices: a Yocto-built
platform image hosting a containerised stack of Mosquitto, a purpose-built
ingestion controller, Eclipse Ditto and MongoDB, exercised by a deterministic
three-device simulator under a pre-registered experimental protocol.
\todo{pending data-v1 --- complete the summary with what the evidence actually
demonstrated, after Chapter~\ref{ch:evaluation} is populated.}

\section{Contributions Demonstrated}
\label{sec:concl-contributions}

\noindent \todo{pending data-v1 --- restate the designed contributions of
Section~\ref{sec:intro-contributions} strictly in terms of what the evidence
in Chapter~\ref{ch:evaluation} supports; each demonstrated contribution must
cite its evidence via the claim--evidence matrix. Contributions not supported
by evidence are reported as limitations instead.}

\section{Answers to the Research Questions}
\label{sec:concl-rqs}

\noindent \todo{pending data-v1 --- one concise paragraph per research
question (RQ1--RQ3), consistent with Section~\ref{sec:eval-answers}; no new
claims may appear here that are not in Chapter~\ref{ch:evaluation}.}

\section{Limitations}
\label{sec:concl-limitations}

\noindent The following limitations are inherent to the study design and hold
regardless of the campaign outcomes:

\begin{itemize}
  \item All telemetry is synthetic; no physical wearables, real users or
        clinical-grade signals were involved.
  \item Performance evidence comes from a single native ARM64 cloud
        \ac{VM} with shared vCPUs; no physical single-board or embedded
        hardware was evaluated, and platform evidence from emulation is
        functional only.
  \item The evaluated stack is a minimal digital-twin core; Ditto search and
        connectivity services, horizontal scaling and multi-gateway
        federation were not studied.
  \item The security configuration is a baseline (\ac{TLS}, authentication,
        least-privilege \ac{ACL}s) and was not subjected to a security
        evaluation.
  \item The identity, ledger and marketplace aspects of the broader C2DTA
        vision were out of scope; nothing in this work validates them.
  \item \todo{pending data-v1 --- add limitations that emerge from the
        campaign itself (e.g.\ conditions removed, variance observed,
        exclusions applied).}
\end{itemize}

\section{Future Work}
\label{sec:concl-future}

\noindent Several directions follow naturally from this work:

\begin{itemize}
  \item \textbf{Physical edge hardware.} Repeating the campaign on physical
        ARM64 single-board hardware, for which the platform recipe is
        prepared but which was not validated here.
  \item \textbf{Identity integration.} A minimal decentralised-identity agent
        alongside the twin stack --- kept strictly conditional and out of the
        core in this project --- evaluated first for deployment feasibility
        and resource footprint on ARM64.
  \item \textbf{Richer device ecosystems.} More device types, device
        onboarding/offboarding dynamics, and multi-gateway topologies.
  \item \textbf{Stack breadth.} Enabling and measuring Ditto search and
        connectivity services, and comparing alternative brokers or twin
        frameworks under the same protocol and harness.
  \item \textbf{Longer horizons.} Multi-day soaks, upgrade-in-place of the
        Yocto image, and lifecycle management of digest-pinned services.
\end{itemize}
